\section{Conceptual Framework}
\TemplateTodo{Explain the relationship between the problem, theory, method, data, and research outputs. A conceptual diagram can help this section.}

\begin{figure}[H]
\centering
\includegraphics[width=0.75\linewidth]{example-figure.pdf}
\caption{Research conceptual framework}
\label{fig:framework}
\end{figure}

Figure~\ref{fig:framework} shows the conceptual framework of the study. Review data from the Google Play Store will go through preprocessing, then feature extraction using TF-IDF, and classification using the Na\"ive Bayes Classifier. The classification results are evaluated using a confusion matrix to obtain accuracy, precision, recall, and F1-score values. Example metric results can be seen in Table~\ref{tab:csv-example}, which shows a comparison of metric values before and after system implementation.

\subsection{Example Table from CSV Data}

The template can display tables populated from CSV files. Store the data in the \texttt{references/} folder and use the \texttt{csvsimple} package.

\begin{table}[H]
\centering
\caption{Example table from CSV data}
\label{tab:csv-example}
\csvautotabular{references/example.csv}
\end{table}

Table~\ref{tab:csv-example} shows sample data from the \texttt{references/example.csv} file. Make sure the CSV column count matches the desired layout.

